% Generated draft: requires mathematical and source review.
\documentclass[11pt]{article}
\usepackage[margin=25mm]{geometry}
\usepackage{fontspec}
\setmainfont{Latin Modern Roman}
\usepackage{amsmath,amssymb,mathtools,unicode-math}
\setmathfont{Latin Modern Math}
\usepackage{enumitem,needspace,xurl,hyperref}
\setlength{\parindent}{0pt}
\setlength{\parskip}{5pt}
\setlength{\emergencystretch}{3em}
\newcommand{\N}{\mathbb{N}}
\newcommand{\Z}{\mathbb{Z}}
\newcommand{\Q}{\mathbb{Q}}
\newcommand{\R}{\mathbb{R}}
\newcommand{\C}{\mathbb{C}}
\newcommand{\F}{\mathbb{F}}
\newcommand{\A}{\mathbb{A}}
\newcommand{\PP}{\mathbb{P}}
\newcommand{\E}{\mathbb{E}}
\newcommand{\eps}{\varepsilon}
\newcommand{\dd}{\,\mathrm{d}}
\begin{document}
\section*{UnsolvedMath ID: 701631}
\textbf{Source problem code:} catalogue entry
\subsection{Definitions and mathematical statement}
For a group \(G\), use the commutator convention \([a,b]=a^{-1}b^{-1}ab\), and define iterated commutators recursively by
\[
[x,{}_1 g]=[x,g],\qquad [x,{}_{n+1}g]=[[x,{}_n g],g]\quad(n\geq 1).
\]
An element \(g\in G\) is \emph{almost Engel} if there exists a finite set \(\mathscr E(g)\subseteq G\) such that
\[
\forall x\in G\;\exists n(x,g)\in\Z_{\geq 1}\;\forall n\in\Z_{\geq 1},\qquad
n\geq n(x,g)\ \Longrightarrow\ [x,{}_n g]\in\mathscr E(g).
\]
The finite set may depend on \(g\), but must work for every \(x\); the threshold may depend on both \(x\) and \(g\). Write
\[
A(G)=\{g\in G:g\text{ is almost Engel}\}.
\]
A \emph{linear group} means a subgroup of \(\operatorname{GL}(m,F)\), the group of invertible \(m\times m\) matrices over a field \(F\), for some positive integer \(m\).

\textbf{Question (P. Shumyatsky, 2018).} For every field \(F\), every positive integer \(m\), and every subgroup \(G\leq\operatorname{GL}(m,F)\), is \(A(G)\) a subgroup of \(G\)?
\subsection{Short English statement}
In an arbitrary linear group, must the almost Engel elements form a subgroup? Here an element is almost Engel when, for each starting element, its sufficiently long iterated commutators eventually stay in a finite set depending only on the almost Engel element.
\subsection{Sources}
\begin{itemize}
\item Local UnsolvedMath registry, supplied source record with permanent ID \(701631\), exact source problem code \texttt{catalogue entry}, title \emph{Kourovka Notebook Problem 19.83}.
\item P. Shumyatsky, Problem 19.83, \emph{The Kourovka Notebook (Unsolved Problems in Group Theory)}. The supplied background identifies the 19th issue (2018) and the 21st issue edition (2026), with identifier \texttt{arXiv:1401.0300v48}. The supplied collection description names E. I. Khukhro and V. D. Mazurov as editors.
\item The supplied registry labels this problem \emph{open}. This is a supplied source classification, not an independently verified assessment of its mathematical or literature status.
\end{itemize}
\end{document}
